\section{Problem Formulation and Threat Model}
\label{sec:problem}

This section formalizes the evidence unit, adversary, and evaluation task used throughout the paper. The current empirical study is intentionally restricted to \appfp observations collected under a verified two-state protocol. The independent external-browser surface is already part of the data architecture, but its attack-effect contribution is treated as a separate research question because the present collection does not yet contain enough complete \browserfp pairs for a defensible comparison.

\subsection{Observation Surfaces}

HybridGuard models a mobile hybrid-app session as a set of logically distinct observation surfaces rather than referring to them collectively as ``three endpoints.'' The distinction matters because each surface is controlled by a different software layer, exposes different APIs, and admits different benign sources of variation.

The first surface is \emph{Android Native}. It contains signals obtained from Android framework and system APIs, including release and build identity, hardware family, physical display properties, memory, sensor availability, battery context, package and installation state, and debugging context. These observations describe the host application and operating environment from the perspective of Native code. The second surface is the \emph{WebView host}. It records properties controlled by the embedded-browser container and the application that instantiates it, such as WebView provider and version, the default user agent, JavaScript-bridge topology, package and installation context, WebView debugging, cleartext policy, and application configuration. Although the WebView host executes in the same application process as Native code, its security and versioning semantics differ from those of the Android framework.

The third surface is the \emph{in-app Web runtime}. It is collected by JavaScript executing inside the FeatureApp WebView and includes browser-visible navigator and platform values, runtime user agent, logical screen and viewport geometry, device-pixel ratio, timezone, Canvas and WebGL outputs, language and plugin/MIME surfaces, resource indicators, and automation-facing attributes. This surface resembles a conventional browser fingerprint, but it is produced inside an application-controlled WebView rather than a standalone browser. The fourth surface is the \emph{external browser}. The same Web probe is launched in an independent system browser and produces \browserfp. It is deliberately modeled as a separate context: similarly named APIs need not return identical values because the WebView and the system browser may use different versions, processes, permissions, storage state, viewport settings, network paths, or customization policies.

The current in-app observation can be written as $x_{\mathrm{app}}=(x_N,x_H,x_W)$, where the three components denote the Native, WebView-host, and in-app Web surfaces. Their frozen dimensions are 84, 26, and 67 signals, respectively. A completed external-browser observation $x_B$ contributes another 67 signals. The derived paired observation $x_{\mathrm{pair}}=(x_{\mathrm{app}},x_B)$ enters the \pairedfp view only when both payloads and their receipts, hashes, collection batches, and pair-provenance records satisfy the frozen contract. When $x_B$ is unavailable, the valid App observation remains an explicit App-only record rather than being padded with artificial Browser values.

\subsection{Field Values, Availability States, and Control Metadata}

For every field $f$, HybridGuard stores a value $v_f$ separately from an availability state $s_f$. The state vocabulary distinguishes an observed value from an unsupported API, permission denial, runtime error, timeout, and a field or relation that is not applicable in the current context. Relation evaluation consumes $(v_f,s_f)$ rather than a value alone, so a parser default or an empty serialization cannot silently convert unavailable evidence into a legitimate zero, false value, or empty string.

Collection and experimental metadata are maintained on a separate control plane. Session identifiers, receipts, batch identifiers, canonical payload hashes, pair identifiers, provider labels, intervention tools, experiment phases, stable-group keys, data splits, and verified labels are used for admission, grouping, and audit, but they are excluded from model-visible evidence. This separation prevents unavailable numeric values from being mistaken for zeroes, prevents a failed Browser launch from becoming an apparently complete 244-field sample, and prevents source names or experiment roles from becoming easier decision shortcuts than the intended environment relationships.

\subsection{Controlled Manipulation Adversary}

The evaluated adversary applies a documented intervention after a baseline collection. Depending on the configuration, it can use Chrome DevTools Protocol controls, Playwright, Puppeteer, Stealth plugins, browser emulation, JavaScript interception, or related runtime mechanisms to modify the WebView or in-app Web surface. The controlled configurations target individual or combined properties such as User-Agent, platform, WebDriver, languages, plugin and MIME surfaces, WebGL vendor and renderer strings, resource values, timezone, and screen or viewport metrics.

This adversary model is intentionally bounded. It does not assume compromise of the Android kernel, successful forgery of hardware-backed attestation, or arbitrary coherent rewriting of every collector output. At the same time, HybridGuard does not equate successful tool execution with an observable attack effect. A comparison enters the current analysis only when independent evidence supports the intervention and direct payload comparison confirms at least one change to a catalogued field. This requirement matters because automation and debugging interfaces have legitimate testing uses, some interventions affect fields outside the collector, and a tool may execute without producing the intended mutation under a particular Android or WebView version.

\subsection{Two-State Evaluation Unit}

Following the advisor's requirement, the paper uses exactly two states, $x^{(b)}\rightarrow x^{(a)}$, where $x^{(b)}$ is the baseline observation and $x^{(a)}$ is collected while the intervention is active. Historical \texttt{clean\_post} records are not used as eligibility conditions, denominators, or evidence for the conclusions in this manuscript. A qualified pair binds the baseline and active observations to the same scenario group and compatible release and schema context while retaining the exact intervention configuration on the control plane.

Each executable relation combines an applicability condition with a comparison and returns one of four states. A \emph{consistent} or \emph{inconsistent} result records the outcome of an applicable comparison, while a \emph{context} result preserves relevant evidence that is not itself a strong alert. The evaluator returns \emph{not assessed} when required evidence is unavailable, unparseable, or outside the declared scope. The principal current observation is a transition from a baseline with no strong relation alert to an active observation with such an alert. That transition describes one frozen relation set on one verified controlled pair; it is not, by itself, a population-level detection metric.

\subsection{Task Boundary and Non-Goals}

The current task is to determine which explicit relationships among the Native, WebView-host, and in-app Web surfaces are violated by verified controlled manipulations, which relations cannot be evaluated, and which confirmed field changes remain outside the frozen relation set. It is therefore distinct from identifying a physical device or tracking a user across sites, proving that an account action is fraudulent, detecting every script that performs fingerprinting, reconstructing complete WebView information flows, classifying malware, or providing cryptographic remote attestation. The study also does not assume that any difference between the in-app runtime and the external browser is malicious, and it does not evaluate the end-to-end effectiveness or user cost of a production RBA policy.

RBA remains the motivating downstream setting. A well-validated environment-integrity signal could inform step-up authentication, restricted access, or manual review. Demonstrating that policy benefit would require account-level ground truth, representative benign and malicious traffic, calibrated thresholds, and user-impact measurements that are not available in the present controlled dataset.

\subsection{Research Questions}

The study is organized around four research questions. \emph{RQ1} asks whether fixed \appfp and \browserfp observations can be governed through explicit field states, release locks, receipts, hashes, and provenance-complete pairing while retaining valid App observations when Browser collection fails. \emph{RQ2} asks which cross-layer \appfp relations change under verified baseline-to-active manipulations and which configurations remain uncovered or unassessed. \emph{RQ3} asks what incremental evidence and benign variability the independently collected \browserfp surface introduces beyond \appfp once sufficiently complete paired data are available. Finally, \emph{RQ4} asks how official-derived relations, device-mined predicates, raw-feature baselines, explicit consistency features, and knowledge-grounded reasoning compare when evaluation is separated by stable device or profile and by controlled scenario. The current results answer RQ2 at a mechanism level, the implemented data path partially addresses RQ1, and RQ3--RQ4 require the future grouped study specified later in the paper.
